The distinction between the two middle values is essential here.
Nevertheless observer union does not enlarge this particular family.

\begin{lemma}[Pure realizations of canonical extensions]
\label{core:lem:essential-extension-pure}
For every nonconstant monotone Boolean function $f$, the operation $D_f$
has a term realization using $\wedge,\vee,\star_a$ if and only if it
has one using only $\wedge,\vee$.
\end{lemma}
\begin{proof}
Distribute the two bichain operations over observer union. Every pure
component has Boolean restriction $f$, since a union of endpoint values
is an endpoint only when all those values agree. Every essential marked
coordinate forces $\pv$ in each component: fix the other free
coordinates at a pivotal Boolean edge, then enlarge them to their middle
values in observer-union order. Both operations preserve this order.
Conversely, no component can produce $\pv$ where $D_f$ gives
$\nv$, since $\pv$ absorbs observer union. Stable outputs
are fixed by the cube endpoints. Each pure component therefore equals
$D_f$ on every input.
\end{proof}

The canonical extension is always attainable at each individual input,
even when no single word attains it everywhere.

\begin{proposition}[Boolean cofactor lifting and the canonical envelope]
\label{core:prop:essential-cofactor-lifting}
Fix a nonconstant monotone Boolean function $f$.
Every pure bichain realization of a Boolean cofactor of $f$ extends to
a pure term with Boolean restriction $f$, whose literal cofactor is
the chosen realization. Endpoint constants are permitted in the cofactor
word. Consequently, for every bichain input $v$
there is a pure term $t$ with Boolean restriction $f$ such that
$t(v)=D_f(v)$.
\end{proposition}
\begin{proof}
For one fixed coordinate $x$, write $f_0\leq f_1$ for the two Boolean
cofactors. If $s$ realizes $f_0$ and $t_1$ realizes $f_1$, the term
$(x\wedge t_1)\vee s$ computes $f$ and restricts exactly to $s$
at $x=b$. To lift a realization $s$ of $f_1$, use
$(x\vee t_0)\wedge s$, which restricts exactly to $s$ at $x=c$.
Iterate over the fixed coordinates. Endpoint constants simplify away
because the final Boolean function is nonconstant.

For an input $v$ with $D_f(v)=\nv$, restrict its stable
coordinates first. Every coordinate marked $\pv$ is inessential
in this cofactor, so choose a formula using only its essential
coordinates. Those inputs are all $\nv$, and the formula gives
$\nv$. Its lift has the required value at $v$.
At all other inputs every pure word computing $f$ already gives $D_f$,
by stable endpoints and the essential marked-source bound.
\end{proof}

An entire marked-source fragment has a realizer when its marked variables
never occur together in a prime minterm. The other variables may have
arbitrary quiet values $b,\nv,c$.

\begin{theorem}[Linear marked-source fragments]
\label{core:thm:linear-marked-fragments}
Partition a finite variable set as $X\sqcup Y$, and let $f$ be a
nonconstant monotone Boolean function. If each prime minterm of $f$
contains at most one coordinate of $X$, there is a pure term $T$ with
\[
 T(\nu)=D_f(\nu)\qquad
 \text{whenever }\{u:\nu_u=\pv\}\subseteq X.
\]
The same conclusion holds if each prime clause contains at most one
coordinate of $X$.
\end{theorem}
\begin{proof}
The minterm hypothesis gives
\[
 f(x,y)=a_0(y)\vee\bigvee_{j\in X}\bigl(x_j\wedge a_j(y)\bigr),
\]
with monotone coefficient functions on $Y$: take $a_0(y)=f(0_X,y)$
and $a_j(y)=f(1_{\{j\}},y)$. For $S\subseteq Y$, define expressions
in decreasing order of $|S|$ by
\begin{equation}\label{core:eq:linear-marked-expansion}
 K_S=
 \begin{cases}
 c,&a_0(S)=1,\\
 \displaystyle\bigvee_{j:\,a_j(S)=1}x_j\ \vee\
 \displaystyle\bigvee_{y\in Y\setminus S}
       \bigl(y\wedge K_{S\cup\{y\}}\bigr),&a_0(S)=0.
 \end{cases}
\end{equation}
Here coefficient values use characteristic assignments, and an empty
disjunction is $b$. Unfolding the finite subset recursion gives a term.
Its Boolean value is $f$: an expansion path adds true parameter variables
until a true baseline coefficient or an available marked coefficient
is reached. Monotonicity makes this equivalent to evaluating $f$.

Fix an input in the stated fragment and its Boolean cube $[L,U]$.
Any marked source reaching the root comes from a leaf $x_j$ in some
$K_S$, with $a_j(S)=1$. Its path uses the parameter gates in $S$, so
survival forces $S\subseteq U\cap Y$. Suppose $j$ is inessential.
Set the parameters to $W=S\cup(L\cap Y)$, retain all fixed true
$X$-coordinates, and set every other free $X$-coordinate except $j$ to
zero. Setting $j=1$ makes $f$ true, so inessentiality makes it true at
$j=0$ too. The displayed linear form then gives either $a_0(W)=1$ or
$a_k(W)=1$ for a fixed true $x_k\in L\cap X$.
Consequently $K_W=c$ on the original input.

Starting at $K_S$, add the parameters in $W\setminus S$ along an
expansion path. Each is fixed $c$, so every gate on that path passes
$K_W=c$ back to $K_S$; an earlier baseline $c$ has the same effect.
Thus $K_S=c$, absorbing the putative marked leaf. No inessential marked
source can reach the root. Stable outputs and essential marked sources
already have their canonical values, so $K_\varnothing=D_f$ on the
whole fragment. Its endpoint constants simplify away because $f$ is
nonconstant. Boolean duality proves the prime-clause version.
\end{proof}

\begin{corollary}[When Boolean data determine the whole operation]
\label{core:cor:essential-unique-fiber}
Let $f$ be nonconstant with full essential support $V$. The bichain term operations
whose Boolean restriction is $f$ form a singleton if and only if
$f(x)=[\sum_{i\in V}x_i\geq r]$ for some $1\leq r\leq|V|$.
This remains true when observer union is allowed. In the singleton
case its unique member is $D_f$.
\end{corollary}
\begin{proof}
For a uniform threshold, every free coordinate of every nonconstant
cube is essential. All pure words computing $f$ therefore give $D_f$;
observer-union expansion gives the same conclusion for full words.

Otherwise $f$ is not constant on some cardinality layer. Connecting
equal-size sets by one-element exchanges gives a set $S$ and distinct
$i,j\notin S$ with $f(S\cup\{i\})=1$ and
$f(S\cup\{j\})=0$. On the two-dimensional cube
$[S,S\cup\{i,j\}]$, its restriction is the projection $i$.
Put $\nv$ at $i$ and $\pv$ at $j$. Cofactor lifting
gives a pure word computing $f$ with value $\nv$ at this input.
Replacing that word $t$ by $t\wedge(t\vee x_j)$ preserves its
Boolean restriction but changes the displayed value to $\pv$.
Thus even the pure fiber is not a singleton.
\end{proof}

\paragraph{A separating source must persist.}
Canonical realizability imposes a condition across different Boolean cubes.
Write $f(A)$ for the value at the characteristic assignment of $A$.

\begin{proposition}[A stable separating source]
\label{core:prop:canonical-separating-source}
If $D_f$ is realized by a bichain term, including observer union, then
for every $f(A)=1$ and $f(B)=0$ there is one $j\in A\setminus B$
essential on every cube $[L,U]$ with $L\subseteq B$, $A\subseteq U$,
and $L\subseteq U$.
\end{proposition}
\begin{proof}
Follow a separating path in a realizing term. At a join choose a child
true at $A$; both children are false at $B$. At a meet choose a child
false at $B$; both are true at $A$. An observer union giving endpoints
at both assignments has separating children, so choose either one.
The resulting leaf is a variable $j\in A\setminus B$, chosen before
the cube.

Mark $j$ on any displayed cube and leave its other free coordinates
neutral. The relation $x=c\Rightarrow y\ne b$ is preserved
coordinatewise by all three bichain operations. Along the chosen path, a meet sibling cannot be the false
endpoint: it is true at $A$, whose true coordinates are not fixed false
on the cube. Dually a join sibling cannot be the true endpoint, since
it is false at $B$ and every coordinate fixed true lies in $B$.
The marked value therefore survives to the root, also through observer
union. Canonical semantics makes $j$ essential on that cube.
\end{proof}

This condition excludes a five-variable operation even though fixing
any one Boolean coordinate gives an exact read-once realizer.

\begin{theorem}[A Boolean-cofactor obstruction]
\label{core:thm:canonical-cofactor-obstruction}
Put
\[
 F(a,b,c,d,z)=ab\lor cd\lor zad\lor zbc.
\]
Its canonical extension $D_F$ is not a bichain term operation, even
with observer union. Each of its ten one-variable Boolean cofactors
has a pure read-once term realizing its canonical extension.
\end{theorem}
\begin{proof}
Take $A=\{a,d,z\}$ and $B=\{b,c\}$. Then $F(A)=1$ and $F(B)=0$.
On the compatible cube with $c=1,b=0$, the function is the projection
$d$; on the compatible cube with $b=1,c=0$, it is the projection $a$.
Their essential sets are disjoint, contradicting
\cref{core:prop:canonical-separating-source}.

The cofactor identities below have read-once right sides:
\[
\begin{array}{c|c}
 a=0&c\wedge(d\vee(z\wedge b))\\
 a=1&b\vee(d\wedge(c\vee z))\\
 z=0&(a\wedge b)\vee(c\wedge d)\\
 z=1&(a\vee c)\wedge(b\vee d).
\end{array}
\]
The permutations $(a\ b)(c\ d)$ and $(a\ c)(b\ d)$ preserve $F$
and act transitively on its first four coordinates, giving the other
six cofactors. A pure read-once formula has exact canonical semantics.
Stable outputs are determined by the cube corners. At a possible output,
choose a child carrying a marked source; its sibling is not a stable
absorber. The disjoint supports let a pivotal edge in that child combine
with a Boolean endpoint assignment for its sibling. Conversely, an
essential marked coordinate forces possible in every Boolean realizer.
Induction therefore gives equality with $D_f$. Apply this to the
displayed cofactors and their permutations.
\end{proof}

The same operation also has exact realizers on every two-source marked
fragment, even with arbitrary neutral inputs elsewhere.

\begin{corollary}[Pairwise marked-source locality]
\label{core:cor:canonical-pair-locality}
For each set $J$ of at most two coordinates of $F$, there is one pure
term $t_J$ such that
\[
 t_J(v)=D_F(v)\qquad
 \text{whenever }\{i:v_i=\pv\}\subseteq J.
\]
\end{corollary}
\begin{proof}
Let $Z_J$ be the semantic budget of $F$ restricted to marks in $J$.
Its direct graph is \eqref{core:eq:direct-marked-edge}. We show that
$F$ respects every direct edge. Its minterms are $ab,cd,zad,zbc$, and
its maximal false facets are $bc,ad,zac,zbd$. Every minterm-minus-facet
has size one, except $zad\setminus bc$ and $zbc\setminus ad$, which
have size three.

If a direct edge crossed from true to false, shrink its source to a
minterm $A$ and enlarge its target to a false facet $B$. The same row
witnesses remain valid, so $A\setminus B\subseteq J$ and this difference
must be a singleton $\{j\}$. But $j$ is essential on every compatible
cube $L\subseteq B$, $U\supseteq A$: the assignment
$L\cup(A\setminus\{j\})\subseteq B$ is false, and adding $j$ makes
it true. No row of $Z_J$ can witness the required inessentiality.

By \cref{core:thm:two-marked-budgets}, $F$ therefore has a pure formula
safe on all of $Z_J$. Its endpoint constants simplify away because $F$
is nonconstant, giving the required $t_J$.
Stable outputs agree with the cube endpoints, and an essential marked
coordinate forces $\pv$ in every Boolean realizer. For the remaining
inputs, a term returns $\pv$ exactly when it does so with just one of
the marked coordinates retained: replacing other $\pv$ inputs by
$\nv$ leaves all stable-endpoint tests unchanged, and each gate
passes the disjunction of its unabsorbed marked-source flags.
Every such single-mark input is in $Z_J$, where the formula is safe.
Thus all inputs in the stated fragment agree with $D_F$.
\end{proof}


\paragraph{Canonical semantics can require repeated sources.}
Exact dependence information does not make a distinguished input linear.
The obstruction is visible in the Boolean contexts in which that input
matters.

Fix a finite passive variable set $V$ and a distinguished input $x\notin V$.
Write $f(\epsilon,A)$ for the Boolean value with distinguished input
$x=\epsilon$ and passive true set $A\subseteq V$, and put
\[
 \mathcal E_x(f)=\{A\subseteq V:f(0,A)=0,\ f(1,A)=1\}.
\]
Define a graph on $\mathcal E_x(f)$ by joining $P,Q$ when, in either
orientation, some passive cube $[L,U]$ satisfies
\begin{equation}\label{core:eq:essential-context-conflict}
 P\subseteq U,\qquad L\subseteq Q,\qquad
 \mathcal E_x(f)\cap[L,U]=\varnothing.
\end{equation}
On such a cube $x$ is inessential. Its common Boolean cofactor is
nonconstant: monotonicity and the two essential contexts force value
$0$ at $L$ and value $1$ at $U$.

\begin{proposition}[Essential-context coloring]
\label{core:prop:essential-context-coloring}
Let $f$ be a nonconstant monotone Boolean function. Every core term whose
four-state restriction is pointwise contained in $D_f$ uses $x$ at least
as many times as the chromatic number of the graph
\eqref{core:eq:essential-context-conflict}. Choices by additional owners
and ordinary union are allowed.
\end{proposition}
\begin{proof}
Choose distinct owners $r_0,r_1,k$. Use endpoint menus
$\{r_0\},\{r_1\}$, neutral $\nv=\{r_0,r_1\}$, and
marked $\pv=\{r_0,r_1,k\}$. Containment in $D_f$
forces the term to compute $f$ on Boolean inputs. At a context
$E\in\mathcal E_x(f)$, mark $x$ and keep the passive inputs Boolean.
The output contains the marker: otherwise omission independence
(\cref{core:prop:omission}) would
let us replace $x$ by either endpoint without changing the value,
contrary to essentiality. Choose an $x$-occurrence supplying this marker
along a surviving path.

One occurrence cannot serve adjacent contexts $P,Q$. At every
$r_0$-owned gate on its path, the sibling is not the absorbing atom $r_0$
at context $P$. At the crossing cube its $r_0$-atomic input coalition is
$V\setminus U\subseteq V\setminus P$; $x$ is marked in both inputs.
Monotonicity of the sibling's singleton profile keeps that flag false.
At an $r_1$-owned gate use context $Q$ instead: its $r_1$-atomic coalition
shrinks from $Q$ to $L$, so that sibling still is not $r_1$.
All other owners act as union on the four-state subalgebra. Hence no
path gate absorbs the marked source on the crossing cube. The marker
survives, contradicting containment in $D_f$, which there equals $\nv$.
Choosing one surviving occurrence at every essential context therefore
gives a proper coloring by the $x$-occurrences.
\end{proof}

\begin{theorem}[Unbounded source repetition in canonical extensions]
\label{core:thm:canonical-source-repetition}
For $m\geq2$, take disjoint variable pairs $\{a_i,b_i\}$ and a new
input $x$. Put
\[
 A_i=a_i\wedge b_i,\qquad
 O_i=\bigvee_{j\ne i}(a_j\vee b_j),\qquad
 T_m=\bigvee_{i=1}^m\bigl(A_i\wedge(x\vee O_i)\bigr),
\]
and let $F_m$ be its Boolean function. Then $T_m$ realizes $D_{F_m}$,
and every realization of $D_{F_m}$ has at least $m$ occurrences of $x$.
The displayed term attains the bound. This holds for every fixed
$K\geq3$, even when all core choices and union are available.
\end{theorem}
\begin{proof}
Each summand is read-once and therefore canonical. To prove that their
outer join is canonical, consider a nonconstant Boolean cube. No summand
is constantly $1$, so the outer join carries a marker exactly when some
summand does. Stable outputs are fixed by the cube endpoints. It therefore
suffices to show that a variable essential for any summand is essential
for $F_m$ on that cube.

If the lower corner already contains a full pair $\{a_h,b_h\}$,
nonconstancy forces it to contain no other variable: $x$ or any outside
bit fixed to $1$ would make the $h$-summand constantly $1$.
The restricted function is consequently the disjunction of $x$ and the
outside variables still free. Every free variable is essential.

Otherwise no pair is already full. Suppose a variable in pair $i$ is
essential for its summand. Set its mate to $1$ and leave all other
pairs incomplete at their lower values. Activate the second factor
by $x$ if available, otherwise by an outside bit already fixed to $1$,
otherwise by one available outside bit with its mate set to $0$.
The last choice is possible because in that case no outside bit was
fixed to $1$. Toggling the queried variable now changes $F_m$ from
$0$ to $1$.

If the essential variable is instead $x$ or an outside variable in the
second factor of summand $i$, no other input of that factor is fixed
to $1$. Set pair $i$ to $1$, all outside variables to $0$ except the
queried variable, and toggle that variable. A single outside bit cannot
complete another pair. Again $F_m$ changes from $0$ to $1$.
These pivotal edges prove canonicality.

The $x$-essential contexts are exactly the $m$ pure pair-sets
$P_i=\{a_i,b_i\}$: with $x=1$ some pair must be full, whereas with
$x=0$ a full pair wins as soon as any outside bit is present.
For $i\ne j$, the cube
\[
 L=\{a_j\},\qquad U=P_i\cup\{a_j\}
\]
contains no $P_h$, but $P_i\subseteq U$ and $L\subseteq P_j$.
The graph of \cref{core:prop:essential-context-coloring} is therefore
the complete graph on $m$ vertices. Its coloring bound gives $m$
occurrences; the displayed construction uses exactly $m$.
\end{proof}


Two-level formulas have a precise matroid boundary. For a finite
nonempty family $\mathcal C$ of nonempty subsets covering $V$, write
\[
 f_{\mathcal C}=\bigwedge_{C\in\mathcal C}\bigvee_{i\in C}x_i,
 \qquad
 T_{\mathcal C}=\bigwedge_{C\in\mathcal C}\bigvee_{i\in C}x_i,
\]
interpreting the first expression in the Boolean algebra and the second
in the bichain. Define its loopless closure operator by
\[
 c_{\mathcal C}(E)=
 V\setminus\bigcup\{C\in\mathcal C:C\cap E=\varnothing\}.
\]
This is the intersection of the clause complements containing $E$.
A matroid closure additionally obeys exchange:
$i\in c(E\cup\{j\})\setminus c(E)$ implies
$j\in c(E\cup\{i\})$.

\begin{theorem}[Matroid boundary for canonical CNFs]
\label{core:thm:essential-cnf-matroid}
The literal CNF $T_{\mathcal C}$ equals $D_{f_{\mathcal C}}$ if and
only if $c_{\mathcal C}$ is a matroid closure. Equivalently, the
inclusion-minimal clauses $\mathcal D$ are the cocircuits of a loopless
matroid $M$, and each clause in $\mathcal C$ is the union of its
contained members of $\mathcal D$. In that case $f_{\mathcal C}$ is
the spanning predicate of $M$.
\end{theorem}
\begin{proof}
First consider the cocircuit CNF of a loopless matroid. It tests spanning,
since a set spans exactly when it meets every cocircuit. On a nonconstant
cube $L\subseteq U$, the set $U$ spans and $L$ does not. No outer
meet operand is therefore $b$. A marked position $j$ produces
$\pv$ exactly when some cocircuit contains $j$ and misses $L$,
equivalently $j\notin\operatorname{cl}_M(L)$.

This is also exactly Boolean essentiality of $j$ on that cube.
Membership in $\operatorname{cl}_M(L)$ prevents $j$ from changing
whether any superset of $L$ spans. Otherwise extend a basis of $L$
together with $j$ to a basis $B\subseteq U$. The set
$L\cup(B\setminus\{j\})$ is nonspanning but becomes spanning after
adding $j$. This is a pivotal edge inside the cube. The CNF thus has
exactly the marked-source observations of $D_{f_{\mathcal C}}$.

We next characterize harmless redundant clauses. A clause $C$ that is
the union of its contained minimal clauses cannot add a marked observation:
if $C$ contains a marked $j$ and misses $L$, some contained minimal clause
also contains $j$ and misses $L$. Stable outputs are unchanged by removing
redundant Boolean clauses. Conversely, suppose $C$ contains a $j$ not
covered by any minimal clause inside $C$. Put $\pv$ at $j$,
$c$ outside $C$, and $\nv$ at all other positions.
Every minimal clause containing $j$ is already satisfied by a fixed $c$.
The restricted Boolean function is nonconstant, because $C$ contains a
minimal clause avoiding $j$, and is independent of $j$. Its canonical
value is $\nv$. Yet the literal clause $C$ gives $\pv$,
and no clause gives $b$, so $T_{\mathcal C}$ gives $\pv$.
Hence equality requires, and allows us to assume, the union condition.

It remains to force cocircuit elimination in the clutter $\mathcal D$.
Suppose distinct $C_1,C_2\in\mathcal D$ contain $j$, but no member
of $\mathcal D$ lies in $(C_1\cup C_2)\setminus\{j\}$.
Choose $i\in C_2\setminus C_1$. Put $\pv$ at $i$, $b$ on
$C_1\setminus\{j\}$, $c$ outside $C_1\cup C_2$, and
$\nv$ at every remaining position. The clause $C_1$ restricts
to the free coordinate $j$. Every clause avoiding $j$ meets the fixed-$c$
region, while every other unsatisfied clause contains $j$. Thus the
Boolean cofactor is exactly $j$, making the marked $i$ inessential.
No clause is constantly $b$, by the antichain property, and $C_2$
produces $\pv$. Again the literal CNF differs from its canonical
extension.

The nonempty antichain $\mathcal D$ must therefore satisfy cocircuit
elimination, so it is a matroid's cocircuit family. Its coverage of $V$
excludes loops. These standard cocircuit axioms and the closure--hyperplane
correspondence are as in \cite[Sections~1.1 and~2.1]{Oxley2011}.
Finally, a set is a union of contained cocircuits precisely when its
complement is a flat: $j\notin\operatorname{cl}_M(E)$ is equivalent
to a cocircuit containing $j$ and missing $E$. Thus the displayed closure
is the matroid closure under the two stated clause conditions.
Conversely, if $c_{\mathcal C}$ is a matroid closure, all clause
complements are flats; its hyperplanes, as maximal proper intersections
of these complements, are exactly the complements of minimal clauses.
This recovers both conditions.
The first part and the redundant-clause argument prove sufficiency.
\end{proof}

\begin{corollary}[Both prime normal forms]
\label{core:cor:essential-both-normal-forms}
Let $f$ be a nonconstant monotone Boolean function on its full essential
support $V$. Its prime CNF and prime DNF both realize $D_f$ if and only
if $f(x)=[\sum_{i\in V}x_i\geq r]$ for some $1\leq r\leq|V|$.
\end{corollary}
\begin{proof}
If the prime CNF is canonical, the theorem identifies $f$ with spanning
in a loopless matroid $M$ of rank $r$. If $M$ is nonuniform, take a
dependent $r$-set and a circuit $C$ inside it. For $j\in C$, put
$\pv$ at $j$, $c$ on $L=C\setminus\{j\}$ and $\nv$
elsewhere. The set $L$ is nonspanning and $j\in\operatorname{cl}_M(L)$,
so $D_f=\nv$. But a basis containing $j$ supplies $\pv$
in the prime DNF. No basis lies in $L$, so no monomial is the absorbing
$c$ and the DNF gives $\pv$. Thus $M$ must be uniform.
Conversely, the prime CNF of a uniform threshold is its uniform matroid's
cocircuit formula. Interchanging $b,c$ and $\wedge,\vee$ turns its prime
DNF into the same construction for the dual threshold of rank
$|V|-r+1$. Both are therefore canonical.
\end{proof}

Disjoint Boolean modules can also be separated without losing any
marked-source information.

\begin{proposition}[Disjoint substitution]
\label{core:prop:essential-disjoint-substitution}
Let $f$ and $g_1,\ldots,g_m$ be nonconstant monotone Boolean functions,
with the $g_i$ on pairwise disjoint variable sets. Then
\[
 D_{f(g_1,\ldots,g_m)}=D_f(D_{g_1},\ldots,D_{g_m}).
\]
If every coordinate of $f$ is essential, the composite canonical
extension is a bichain term operation if and only if $D_f$ and every
$D_{g_i}$ are bichain term operations. The same equivalence holds
with observer union allowed.
\end{proposition}
\begin{proof}
Each block cube induces a fixed Boolean bit or a free outer coordinate.
Disjointness lets these coordinates range independently over the entire
outer cube. A variable in block $i$ is essential for the composite
exactly when it is essential in $g_i$ and coordinate $i$ is essential
for $f$ on the induced cube. Indeed a pivotal composite edge supplies
both pivotal edges; conversely combine the two pivotal edges and
independently realize the other block outputs. This proves the displayed
identity, including stable outputs, and hence sufficiency by composition.

For necessity, fix all but a pivotal variable in each block at Boolean
endpoints. This recovers $f$ as a cofactor of the composite. To recover
$g_i$, fix the other blocks to endpoint outputs along a pivotal edge
for outer coordinate $i$. Canonical extension commutes with Boolean
cofactors. In a pure term, endpoint constants simplify away whenever
the resulting Boolean function is nonconstant. The pure-realization
lemma handles observer union as well.
\end{proof}

